% ch28_formal_yoneda.tex — Volume III Chapter 4

\chapter{The Formal Yoneda Lemma}

\begin{overviewbox}
Yoneda has appeared three times: as a hom-set bijection (Chapter~9), as a
representability principle (Chapter~15), and in enriched form (Chapter~20).
These three statements are facets of one theorem. With ends and coends in
hand (Chapter~26) and weighted (co)limits available (Chapter~27), we can
unify them.

The \emph{formal} Yoneda lemma states a single isomorphism in any sufficiently
complete enriched setting:
\[
  [\mathcal{C}, \mathcal{V}]_\mathcal{V}(\mathcal{C}(a, -), F)
  \;=\; \int_c \mathcal{V}(\mathcal{C}(a, c), F c)
  \;\cong\; F(a).
\]
The proof is now mechanical: apply the ninja Yoneda lemma. From this we
derive: the \term{density theorem} (every presheaf is a coend of representables);
the \term{free cocompletion} property of the presheaf category; the
\term{coend Yoneda} and \term{end Yoneda} forms; and the \emph{ninja
duality} between left and right Kan extensions through the Yoneda embedding.
The chapter closes with the formal Yoneda lemma in the 2-categorical
language of $\mathbf{Prof}$.
\end{overviewbox}


% ═══════════════════════════════════════════════════════════════════════════
\begin{nameorigin}
The \term{ninja Yoneda lemma} is a popular informal name for the coend
form of Yoneda: $\int^c \mathcal{C}(c, c') \cdot F(c) \cong F(c')$.
The name ``ninja'' echoes the way the lemma is invoked in proofs:
suddenly, without warning, a coend collapses to a single value — like
a ninja striking. The terminology is jocular but universal in
category-theoretic blog posts, lectures, and even some textbooks. The
formal name is \term{co-Yoneda} or \emph{Yoneda reduction}.

\term{Density theorem} (cross-ref Vol II Ch.~24, Vol IV Ch.~48) is the
coend / colimit statement: every presheaf is a canonical colimit of
representables. The terminology echoes topology (a ``dense'' subset has
closure equal to the whole space); categorically, representables are
dense in $\mathcal{P}(\mathcal{C})$.
\end{nameorigin}

\section{The Unified Statement}
\label{sec:unified-yoneda}
% ═══════════════════════════════════════════════════════════════════════════

\subsection{Formalizing}

\begin{theorem}[Yoneda lemma, unified form]
\label{thm:yoneda-unified}
Let $\mathcal{V}$ be a complete symmetric monoidal closed category and
$\mathcal{C}$ a small $\mathcal{V}$-category. For any $\mathcal{V}$-functor
$F : \mathcal{C} \to \mathcal{V}$ and object $a \in \mathcal{C}$, there is a
$\mathcal{V}$-natural isomorphism
\begin{equation*}
  \mathrm{Yo}_F^a \;:\; \int_c \mathcal{V}(\mathcal{C}(a, c), F c)
  \;\xrightarrow{\;\cong\;}\; F(a)
\end{equation*}
in $\mathcal{V}$, natural in $a$ and in $F$.
\end{theorem}

\begin{prooflog}
\proofstep{The end on the left is precisely the $\mathcal{V}$-object of
$\mathcal{V}$-natural transformations: $[\mathcal{C}, \mathcal{V}]_\mathcal{V}
(\mathcal{C}(a, -), F)$.}{End formula for enriched naturals (Chapter~26).}
\proofstep{We construct a candidate map $\mathrm{Yo}_F^a : [\mathcal{C}, \mathcal{V}]_\mathcal{V}
(\mathcal{C}(a, -), F) \to F(a)$ by evaluation at $a$: precompose with the
identity $\mathcal{V}$-morphism $j_a : I \to \mathcal{C}(a, a)$.}{This is the
forward direction; explicit on representables.}
\proofstep{Construct the inverse via the $\mathcal{V}$-functoriality of $F$:
given $x : I \to F(a)$ in $\mathcal{V}$, define $\widetilde x_c :
\mathcal{C}(a, c) \to F(c)$ as the composite $\mathcal{C}(a, c) \cong I \otimes
\mathcal{C}(a, c) \xrightarrow{x \otimes \mathrm{id}} F(a) \otimes \mathcal{C}(a, c)
\xrightarrow{F_{a,c}} F(c)$.}{Uses the $\mathcal{V}$-functor structure
$F_{a,c} : F(a) \otimes \mathcal{C}(a, c) \to F(c)$.}
\proofstep{The family $(\widetilde x_c)_c$ is a wedge to $T(c, c') =
\mathcal{V}(\mathcal{C}(a, c), F c')$ — i.e., it satisfies the dinatural
condition.}{By the $\mathcal{V}$-functor axioms of $F$ applied to morphisms
$c \to c'$ in $\mathcal{C}$.}
\proofstep{So $(\widetilde x_c)$ factors through the end to give a morphism
$\widetilde x : I \to \int_c \mathcal{V}(\mathcal{C}(a, c), F c)$.}{Universal
property of the end.}
\proofstep{$\mathrm{Yo}_F^a \circ \widetilde{(-)} = \mathrm{id}_{F(a)}$:
evaluating $\widetilde x$ at $a$ gives $F_{a,a}(x \otimes j_a) = x$ by the
$\mathcal{V}$-functor unit axiom.}{Unit-of-$F$.}
\proofstep{$\widetilde{(-)} \circ \mathrm{Yo}_F^a = \mathrm{id}$: starting
from a $\mathcal{V}$-natural transformation, evaluating at $a$ then
transporting via $F_{a,-}$ reproduces the original transformation, by the
$\mathcal{V}$-functor composition axiom (the naturality of $\theta$ matches
$F$'s action exactly).}{Composition axiom of $F$.}
\proofstep{Both composites are identities, so $\mathrm{Yo}_F^a$ is an iso
in $\mathcal{V}$.}{Yoneda established at $a$.}
\proofstep{Naturality in $a$ and $F$ is automatic: all constructions
(end, evaluation, $F$'s action) are $\mathcal{V}$-natural in both arguments.}{
Standard.}
\end{prooflog}

\subsection{Reorganizing: The Three Forms Recovered}

\begin{corollary}[Three forms of Yoneda]
\label{cor:three-yonedas}
Theorem~\ref{thm:yoneda-unified} specializes to:
\begin{itemize}
  \item \textbf{Set-Yoneda} (Chapter~9): $\mathcal{V} = \mathbf{Set}$,
        $\mathcal{C}$ ordinary locally small. Then $\mathrm{Nat}
        (\mathcal{C}(a, -), F) \cong F(a)$.
  \item \textbf{Enriched Yoneda} (Chapter~20): $\mathcal{V}$ symmetric
        monoidal closed, $\mathcal{C}$ a $\mathcal{V}$-category. The
        statement above.
  \item \textbf{Coend Yoneda} (the ninja Yoneda): the dual form,
        $\int^c \mathcal{C}(c, a) \otimes F(c) \cong F(a)$, obtained by
        the end-coend duality for hom-objects.
\end{itemize}
\end{corollary}

\begin{remark}
The ``three forms'' are not three different theorems. They are one theorem
expressed in different language: ends (the explicit hom-set / hom-object
form), set-natural transformations (the original statement), and coends (the
``tensor-out'' dualization). Each is the others' shadow.
\end{remark}


% ═══════════════════════════════════════════════════════════════════════════
\section{The Yoneda Embedding}
\label{sec:formal-yoneda-embedding}
% ═══════════════════════════════════════════════════════════════════════════

\subsection{Formalizing}

\begin{definition}[Enriched Yoneda embedding]
\label{def:formal-yoneda-embedding}
For a small $\mathcal{V}$-category $\mathcal{C}$, the \term{Yoneda embedding}
is the $\mathcal{V}$-functor
\begin{equation*}
  Y \;:\; \mathcal{C} \to [\mathcal{C}^{\mathrm{op}}, \mathcal{V}]_\mathcal{V},
  \qquad
  a \mapsto \mathcal{C}(-, a).
\end{equation*}
\end{definition}

\begin{theorem}[Yoneda embedding is fully faithful]
\label{thm:yoneda-ff}
The Yoneda embedding $Y : \mathcal{C} \to [\mathcal{C}^{\mathrm{op}},
\mathcal{V}]_\mathcal{V}$ is a fully faithful $\mathcal{V}$-functor: it
induces $\mathcal{V}$-isomorphisms
\begin{equation*}
  Y_{a, b} \;:\; \mathcal{C}(a, b) \;\xrightarrow{\;\cong\;}\;
  [\mathcal{C}^{\mathrm{op}}, \mathcal{V}]_\mathcal{V}(\mathcal{C}(-, a), \mathcal{C}(-, b)).
\end{equation*}
\end{theorem}

\begin{proof}
By Theorem~\ref{thm:yoneda-unified} applied to $F = \mathcal{C}(-, b) :
\mathcal{C}^{\mathrm{op}} \to \mathcal{V}$ at the object $a$:
\begin{equation*}
  [\mathcal{C}^{\mathrm{op}}, \mathcal{V}]_\mathcal{V}(\mathcal{C}(-, a), \mathcal{C}(-, b))
  \;\cong\; \mathcal{C}(a, b).
\end{equation*}
That this isomorphism is exactly the hom-functor of $Y$ (i.e., $Y_{a,b}$) is
the content of the construction.
\end{proof}

\subsection{Reorganizing: The Universal Property}

\begin{theorem}[Free cocompletion]
\label{thm:free-cocompletion}
For any cocomplete $\mathcal{V}$-category $\mathcal{E}$ and $\mathcal{V}$-functor
$F : \mathcal{C} \to \mathcal{E}$, there is an essentially unique
colimit-preserving (i.e., $\mathcal{V}$-cocontinuous) $\mathcal{V}$-functor
\begin{equation*}
  \widetilde F \;:\; [\mathcal{C}^{\mathrm{op}}, \mathcal{V}]_\mathcal{V} \to \mathcal{E}
\end{equation*}
with $\widetilde F \circ Y \cong F$. Explicitly, on a presheaf $P$:
\begin{equation*}
  \widetilde F (P) \;=\; P \star F \;=\; \int^c P(c) \otimes F(c).
\end{equation*}
\end{theorem}

\begin{prooflog}
\proofstep{Define $\widetilde F (P) = \int^c P(c) \otimes F(c)$ (a weighted
colimit with $P$ as weight).}{Coend formula.}
\proofstep{$\widetilde F$ is a $\mathcal{V}$-functor by functoriality of
weighted colimits in the weight (Chapter~27).}{Functoriality.}
\proofstep{$\widetilde F \circ Y(a) = \widetilde F(\mathcal{C}(-, a)) =
\int^c \mathcal{C}(c, a) \otimes F(c) \cong F(a)$.}{By the coend Yoneda lemma
(Theorem~\ref{thm:yoneda-unified} dual / ninja Yoneda).}
\proofstep{$\widetilde F$ preserves all small colimits because coends commute
with colimits (in their non-integration arguments).}{Coend = colimit, and
colimits commute with colimits.}
\proofstep{For uniqueness, suppose $G : [\mathcal{C}^{\mathrm{op}}, \mathcal{V}]
_\mathcal{V} \to \mathcal{E}$ is colimit-preserving with $G \circ Y \cong F$.}{
Compare to any candidate.}
\proofstep{Every presheaf is a coend of representables: $P \cong \int^c P(c)
\otimes Y(c)$ (co-Yoneda).}{Density theorem in coend form.}
\proofstep{Applying $G$: $G(P) \cong G\!\left(\int^c P(c) \otimes Y(c)\right)
\cong \int^c P(c) \otimes G(Y(c)) \cong \int^c P(c) \otimes F(c) = \widetilde F(P)$.}{
$G$ preserves colimits; $G \circ Y = F$.}
\proofstep{Therefore $G \cong \widetilde F$, essentially uniquely.}{Uniqueness.}
\end{prooflog}

\begin{corollary}[{Universal property of $[\mathcal{C}^{\mathrm{op}}, \mathcal{V}]_\mathcal{V}$}]
\label{cor:presheaf-universal}
The Yoneda embedding $Y : \mathcal{C} \to [\mathcal{C}^{\mathrm{op}},
\mathcal{V}]_\mathcal{V}$ exhibits the presheaf category as the free
$\mathcal{V}$-cocompletion of $\mathcal{C}$: it is initial among $\mathcal{V}$-functors
from $\mathcal{C}$ into cocomplete $\mathcal{V}$-categories.
\end{corollary}


% ═══════════════════════════════════════════════════════════════════════════
\section{The Density Theorem}
\label{sec:density}
% ═══════════════════════════════════════════════════════════════════════════

\subsection{Formalizing}

\begin{theorem}[Density / co-Yoneda]
\label{thm:density-formal}
For any $\mathcal{V}$-functor $P : \mathcal{C}^{\mathrm{op}} \to \mathcal{V}$,
\begin{equation*}
  P \;\cong\; \int^c P(c) \otimes \mathcal{C}(-, c)
\end{equation*}
as $\mathcal{V}$-functors. Every presheaf is canonically a coend of
representables.
\end{theorem}

\begin{proof}
At each $a \in \mathcal{C}$, the right-hand side evaluates to
$\int^c P(c) \otimes \mathcal{C}(a, c) \cong P(a)$ by ninja Yoneda
(Chapter~26, Theorem~\ref{thm:yoneda-reduction-coend}). Naturality in $a$ is
automatic.
\end{proof}

\subsection{Reorganizing: The Yoneda Embedding is Dense}

\begin{definition}[Dense functor]
\label{def:dense}
A $\mathcal{V}$-functor $K : \mathcal{C} \to \mathcal{D}$ is \term{dense} if
the canonical map
\begin{equation*}
  \int^c \mathcal{D}(K c, -) \otimes K c \;\longrightarrow\; \mathrm{Id}_\mathcal{D}
\end{equation*}
is an isomorphism — i.e., the identity functor of $\mathcal{D}$ is the left
Kan extension of $K$ along itself.
\end{definition}

\begin{corollary}
\label{cor:yoneda-dense}
The Yoneda embedding $Y : \mathcal{C} \to [\mathcal{C}^{\mathrm{op}},
\mathcal{V}]_\mathcal{V}$ is dense.
\end{corollary}

\begin{proof}
The density condition at a presheaf $P$ becomes
\begin{equation*}
  \int^c [\mathcal{C}^{\mathrm{op}}, \mathcal{V}]_\mathcal{V}(\mathcal{C}(-, c), P)
  \otimes \mathcal{C}(-, c) \;\cong\; P.
\end{equation*}
By the Yoneda lemma applied to the inner $[\mathcal{C}^{\mathrm{op}},
\mathcal{V}]_\mathcal{V}(\mathcal{C}(-, c), P) \cong P(c)$, this reduces to
$\int^c P(c) \otimes \mathcal{C}(-, c) \cong P$ — which is the density
theorem.
\end{proof}


% ═══════════════════════════════════════════════════════════════════════════
\section{Adjoint Functors via Yoneda}
\label{sec:adjoints-via-yoneda}
% ═══════════════════════════════════════════════════════════════════════════

\subsection{Formalizing}

The Yoneda lemma gives a clean formulation of when a functor has an adjoint.

\begin{theorem}[Adjoint as represented hom-functor]
\label{thm:adjoint-as-rep}
A $\mathcal{V}$-functor $G : \mathcal{D} \to \mathcal{C}$ has a left adjoint
iff for every $c \in \mathcal{C}$, the $\mathcal{V}$-functor
\begin{equation*}
  \mathcal{C}(c, G-) \;:\; \mathcal{D} \to \mathcal{V}
\end{equation*}
is representable. In that case, the representing object is $L(c)$, where
$L \dashv G$.
\end{theorem}

\begin{proof}
\emph{($\Rightarrow$):} If $L \dashv G$, then $\mathcal{C}(c, G-) \cong
\mathcal{D}(L c, -)$, manifestly representable by $L c$. \\
\emph{($\Leftarrow$):} If $\mathcal{C}(c, G-)$ is represented by $L c$, the
correspondence $c \mapsto L c$ is functorial (by uniqueness of representing
objects), and the isomorphism is exactly the adjoint hom-bijection.
\end{proof}

\begin{corollary}[Right Kan extension formula for adjoints]
\label{cor:adjoint-as-ran}
A right adjoint $G : \mathcal{D} \to \mathcal{C}$ is the right Kan extension
of itself along the Yoneda embedding:
\begin{equation*}
  G \;\cong\; \operatorname{Ran}_Y (G \circ Y^{-1}) \quad\text{informally},
\end{equation*}
or, more precisely: every right adjoint is, up to natural isomorphism,
$\operatorname{Ran}_Y F$ for some $F : \mathcal{C} \to \mathcal{C}$, where
$Y$ is the Yoneda embedding of an appropriate source.
\end{corollary}

\subsection{Reorganizing: Universal Constructions as Kan Extensions}

The combination of Theorem~\ref{thm:adjoint-as-rep}, the free cocompletion
Theorem~\ref{thm:free-cocompletion}, and the weighted Kan extension formulas
(Chapter~27) yields the slogan:
\begin{quote}
\emph{Every universal construction is the left Kan extension of an evident
functor along the Yoneda embedding.}
\end{quote}

\begin{example_thm}[Free abelian group]
The free abelian group functor $\mathbf{Set} \to \mathbf{Ab}$ is
$\operatorname{Lan}_Y (\mathbb{Z}[-])$ where $\mathbb{Z}[-]$ sends a finite
set to the free abelian group on it and $Y$ is the Yoneda embedding of
finite sets into $\mathbf{Set}$. Computing the coend formula recovers the
standard construction.
\end{example_thm}

\begin{example_thm}[Geometric realization]
$|-| : \mathbf{sSet} \to \mathbf{Top}$ is $\operatorname{Lan}_Y(\Delta^\bullet)$
where $\Delta^\bullet : \Delta \to \mathbf{Top}$ sends $[n]$ to the topological
$n$-simplex, and $Y : \Delta \to \mathbf{sSet}$ is the Yoneda embedding
$[n] \mapsto \Delta[n]$. This recovers the coend formula $|X| = \int^{[n]}
X_n \cdot \Delta^n$ from Chapter~26.
\end{example_thm}


% ═══════════════════════════════════════════════════════════════════════════
\section{Yoneda in $\mathbf{Prof}$}
\label{sec:yoneda-prof}
% ═══════════════════════════════════════════════════════════════════════════

\subsection{Formalizing}

The bicategory $\mathbf{Prof}$ (Chapter~25, \S~5.2) houses profunctors as
1-cells. Its unit at $\mathcal{C}$ is the hom-profunctor $\mathcal{C}(-, -) :
\mathcal{C}^{\mathrm{op}} \times \mathcal{C} \to \mathbf{Set}$.

\begin{theorem}[Yoneda lemma in $\mathbf{Prof}$]
\label{thm:yoneda-prof}
For a small category $\mathcal{C}$, the hom-profunctor $\mathcal{C}(-, -) :
\mathcal{C} \to \mathcal{C}$ is the \emph{identity 1-cell} of $\mathbf{Prof}$
at $\mathcal{C}$, in the sense that for any profunctor $P : \mathcal{C}
\to \mathcal{D}$:
\begin{equation*}
  P \circ \mathcal{C}(-, -) \;\cong\; P, \qquad \mathcal{D}(-, -) \circ P \;\cong\; P.
\end{equation*}
\end{theorem}

\begin{proof}
The left identity: $(P \circ \mathcal{C}(-,-))(d, c) = \int^{c'} P(d, c')
\times \mathcal{C}(c', c) \cong P(d, c)$ by the ninja Yoneda lemma applied
to $P(d, -)$. The right identity is the dual computation.
\end{proof}

\begin{remark}[Yoneda is unitality]
This is the most distilled form of Yoneda's content: \emph{the hom-functor
is the identity of profunctor composition}. All the other forms of Yoneda
follow from this single fact, because profunctor composition is the coend
and the unit is the hom-profunctor.
\end{remark}

\subsection{Reorganizing: Cauchy Completeness}

\begin{definition}[Cauchy completeness]
\label{def:cauchy-complete}
A small $\mathcal{V}$-category $\mathcal{C}$ is \term{Cauchy complete} (or
\emph{Karoubian} in the $\mathbf{Set}$-case) if every $\mathcal{V}$-profunctor
$P : \mathcal{C} \to \mathbf{1}$ representable in $\mathbf{Prof}$ is
representable by an object of $\mathcal{C}$ — equivalently, every \emph{absolute}
weighted colimit exists in $\mathcal{C}$.
\end{definition}

\begin{remark}
Cauchy completion is the universal Cauchy-complete enlargement of $\mathcal{C}$.
It is the most refined version of the ``presheaf category as free cocompletion''
theme: not all colimits, but exactly the colimits that are forced by
profunctor representability.
\end{remark}


% ═══════════════════════════════════════════════════════════════════════════
\section{Exercises}
\label{sec:formal-yoneda-exercises}
% ═══════════════════════════════════════════════════════════════════════════

\begin{exercisesbox}
\begin{qa}
\qaitem{State the unified Yoneda lemma.}{For $F : \mathcal{C} \to \mathcal{V}$ a $\mathcal{V}$-functor and $a \in \mathcal{C}$: $\int_c \mathcal{V}(\mathcal{C}(a, c), F c) \cong F(a)$ in $\mathcal{V}$, $\mathcal{V}$-naturally in $a$ and $F$.}
\qaitem{What is the ``end form'' of natural transformations?}{$[\mathcal{C}, \mathcal{V}]_\mathcal{V}(G, H) = \int_c \mathcal{V}(G c, H c)$.}
\qaitem{What is the ``coend form'' of Yoneda (ninja Yoneda)?}{$\int^c \mathcal{C}(c, a) \otimes F(c) \cong F(a)$.}
\qaitem{How does ordinary Set-Yoneda emerge from the unified statement?}{Take $\mathcal{V} = \mathbf{Set}$, $\mathcal{C}$ locally small. The end becomes $\mathrm{Nat}(\mathcal{C}(a, -), F)$ and the conclusion is $\mathrm{Nat}(\mathcal{C}(a, -), F) \cong F(a)$.}
\qaitem{Define the (enriched) Yoneda embedding.}{$Y : \mathcal{C} \to [\mathcal{C}^{\mathrm{op}}, \mathcal{V}]_\mathcal{V}$, $a \mapsto \mathcal{C}(-, a)$.}
\qaitem{Why is the Yoneda embedding fully faithful?}{Because $[\mathcal{C}^{\mathrm{op}}, \mathcal{V}]_\mathcal{V}(\mathcal{C}(-, a), \mathcal{C}(-, b)) \cong \mathcal{C}(a, b)$ by the unified Yoneda lemma applied to $F = \mathcal{C}(-, b)$.}
\qaitem{State the free cocompletion theorem.}{For any cocomplete $\mathcal{V}$-category $\mathcal{E}$ and $F : \mathcal{C} \to \mathcal{E}$, there is an essentially unique colimit-preserving $\widetilde F : [\mathcal{C}^{\mathrm{op}}, \mathcal{V}]_\mathcal{V} \to \mathcal{E}$ with $\widetilde F \circ Y \cong F$.}
\qaitem{What is the explicit formula for $\widetilde F$ on a presheaf $P$?}{$\widetilde F(P) = \int^c P(c) \otimes F(c) = P \star F$.}
\qaitem{State the density theorem in coend form.}{$P \cong \int^c P(c) \otimes \mathcal{C}(-, c)$. Every presheaf is canonically a coend of representables.}
\qaitem{What does ``dense functor'' mean?}{A $\mathcal{V}$-functor $K : \mathcal{C} \to \mathcal{D}$ such that the canonical map $\int^c \mathcal{D}(K c, -) \otimes K c \to \mathrm{Id}_\mathcal{D}$ is an iso — equivalently, $\mathrm{Id}_\mathcal{D} = \operatorname{Lan}_K K$.}
\qaitem{Why is the Yoneda embedding dense?}{Because the density condition reduces, via the Yoneda lemma, to the density theorem itself.}
\qaitem{When does $G : \mathcal{D} \to \mathcal{C}$ have a left adjoint?}{Iff $\mathcal{C}(c, G-) : \mathcal{D} \to \mathcal{V}$ is representable for every $c \in \mathcal{C}$.}
\qaitem{What is the relationship between universal constructions and Kan extensions?}{Every universal construction is a left Kan extension along the Yoneda embedding of an appropriate sub-category.}
\qaitem{Express the free abelian group functor as a Kan extension.}{$F_{\mathrm{ab}} = \operatorname{Lan}_Y (\mathbb{Z}[-])$, where $Y : \mathbf{FinSet} \to \mathbf{Set}$ and $\mathbb{Z}[-] : \mathbf{FinSet} \to \mathbf{Ab}$.}
\qaitem{Express geometric realization as a Kan extension.}{$|-| = \operatorname{Lan}_Y(\Delta^\bullet)$, $Y : \Delta \to \mathbf{sSet}$, $\Delta^\bullet : \Delta \to \mathbf{Top}$.}
\qaitem{What is the identity 1-cell of $\mathbf{Prof}$ at $\mathcal{C}$?}{The hom-profunctor $\mathcal{C}(-, -)$.}
\qaitem{State Yoneda as the unitality of $\mathbf{Prof}$.}{For any profunctor $P : \mathcal{C} \to \mathcal{D}$: $P \circ \mathcal{C}(-,-) \cong P$ and $\mathcal{D}(-,-) \circ P \cong P$.}
\qaitem{Why is ``Yoneda is unitality'' the most distilled form?}{Because all other Yoneda statements (representable functor, hom-bijection, density) follow from this single profunctor-composition fact.}
\qaitem{What is a Cauchy complete category?}{A category in which every absolute weighted colimit exists; equivalently, every profunctor to the terminal category that is representable in $\mathbf{Prof}$ is representable by an object.}
\qaitem{What is the Karoubian (idempotent splitting) view of Cauchy completion?}{In the $\mathbf{Set}$-enriched case, Cauchy completion is the splitting of all idempotents.}
\qaitem{Why is Cauchy completion less than full cocompletion?}{Because only \emph{absolute} colimits are forced; the presheaf category (full cocompletion) typically contains many more objects than the Cauchy completion.}
\qaitem{What is the relationship between $\widetilde F$ and Lan$_Y F$?}{They are equal: $\widetilde F = \operatorname{Lan}_Y F$. The free cocompletion is the left Kan extension along Yoneda.}
\qaitem{Why does the colimit-preserving extension $\widetilde F$ exist and is unique?}{Existence: the coend formula constructs it. Uniqueness: every presheaf is a coend of representables (density), so any colimit-preserving extension is determined on representables and is then forced by coend preservation.}
\qaitem{What does ``natural in $F$'' mean in the unified Yoneda statement?}{The isomorphism $\int_c \mathcal{V}(\mathcal{C}(a, c), F c) \cong F(a)$ is a natural transformation between two $\mathcal{V}$-functors $[\mathcal{C}, \mathcal{V}]_\mathcal{V} \to \mathcal{V}$: the end of the hom and evaluation at $a$.}
\qaitem{What does the Yoneda embedding ``embed'' a category into?}{The presheaf category $[\mathcal{C}^{\mathrm{op}}, \mathcal{V}]_\mathcal{V}$, by viewing each object $a$ as the representable presheaf $\mathcal{C}(-, a)$.}
\qaitem{Why is the Yoneda embedding a $\mathcal{V}$-functor, not just an ordinary functor?}{Because the action on hom-objects is the $\mathcal{V}$-natural enrichment of the action on objects; the $\mathcal{V}$-structure is preserved.}
\qaitem{What is the dual of the free cocompletion theorem?}{The free completion theorem: for any complete $\mathcal{V}$-category $\mathcal{E}$ and $F : \mathcal{C}^{\mathrm{op}} \to \mathcal{E}$, the right Kan extension $\operatorname{Ran}_{Y^{\mathrm{op}}} F : [\mathcal{C}, \mathcal{V}]_\mathcal{V} \to \mathcal{E}$ is the essentially unique limit-preserving extension.}
\qaitem{In what sense is Yoneda ``trivial''?}{Once expressed as the unitality of $\mathbf{Prof}$ or as the ninja Yoneda formula via coend, the proof is mechanical: substitute the identity, apply functoriality, conclude.}
\qaitem{What is the ``coherence'' the unified statement encodes?}{The compatibility of three pictures: hom-natural transformations as ends; representability as Yoneda; density as the co-Yoneda formula. All are facets of a single coend/end identity.}
\qaitem{Why is this chapter called ``Formal Yoneda'' rather than just ``Yoneda Revisited''?}{Because the statement is now formal in the sense of formal category theory: it works in any enriched / 2-categorical setting where ends, coends, and weighted (co)limits make sense.}
\end{qa}
\end{exercisesbox}


% ═══════════════════════════════════════════════════════════════════════════
\section*{Chapter Summary}
\addcontentsline{toc}{section}{Chapter Summary}
% ═══════════════════════════════════════════════════════════════════════════

\begin{summarybox}
\textbf{The unified Yoneda lemma.}\quad For a $\mathcal{V}$-functor $F :
\mathcal{C} \to \mathcal{V}$ and $a \in \mathcal{C}$:
\[
  \int_c \mathcal{V}(\mathcal{C}(a, c), F c) \;\cong\; F(a)
\]
in $\mathcal{V}$. This unifies the set-, enriched-, and coend-forms of
Yoneda into one statement.

\textbf{The Yoneda embedding.}\quad $Y : \mathcal{C} \to [\mathcal{C}^{\mathrm{op}},
\mathcal{V}]_\mathcal{V}$ is fully faithful: the Yoneda lemma identifies
$\mathcal{C}(a, b)$ with $\mathcal{V}$-natural transformations
$\mathcal{C}(-, a) \Rightarrow \mathcal{C}(-, b)$.

\textbf{Free cocompletion.}\quad $[\mathcal{C}^{\mathrm{op}}, \mathcal{V}]_\mathcal{V}$
is the free $\mathcal{V}$-cocompletion of $\mathcal{C}$: every
$\mathcal{V}$-functor $F : \mathcal{C} \to \mathcal{E}$ to a cocomplete
$\mathcal{V}$-category extends uniquely (up to iso) to a colimit-preserving
$\widetilde F = \operatorname{Lan}_Y F$, given by the coend formula
$\widetilde F(P) = \int^c P(c) \otimes F(c)$.

\textbf{Density.}\quad Every presheaf is a coend of representables: $P \cong
\int^c P(c) \otimes \mathcal{C}(-, c)$. The Yoneda embedding is dense: the
identity functor of the presheaf category is the left Kan extension of $Y$
along itself.

\textbf{Adjoints via representability.}\quad A $\mathcal{V}$-functor $G$ has
a left adjoint iff $\mathcal{C}(c, G-)$ is representable for every $c$. Every
adjoint construction is a Kan extension along Yoneda.

\textbf{Yoneda in $\mathbf{Prof}$.}\quad The hom-profunctor is the identity
1-cell of $\mathbf{Prof}$ at $\mathcal{C}$. This is the most distilled form
of Yoneda's content; all other forms follow.

\textbf{Cauchy completeness.}\quad The category where exactly the absolute
weighted colimits exist; in $\mathbf{Set}$-enriched, the splitting of
idempotents.
\end{summarybox}


% ═══════════════════════════════════════════════════════════════════════════
\section*{Formal Restatement}
\addcontentsline{toc}{section}{Formal Restatement}
% ═══════════════════════════════════════════════════════════════════════════

\begin{formalrestatement}
\textbf{Unified Yoneda.}\quad For $F : \mathcal{C} \to \mathcal{V}$ a
$\mathcal{V}$-functor:
\[
  \int_c \mathcal{V}(\mathcal{C}(a, c), F c) \;\cong\; F(a)
  \qquad
  \int^c \mathcal{C}(c, a) \otimes F(c) \;\cong\; F(a)
\]
$\mathcal{V}$-naturally in $a$ and $F$.

\medskip
\textbf{Yoneda embedding.}\quad $Y : \mathcal{C} \to [\mathcal{C}^{\mathrm{op}},
\mathcal{V}]_\mathcal{V}$, $a \mapsto \mathcal{C}(-, a)$. Fully faithful as a
$\mathcal{V}$-functor.

\medskip
\textbf{Free cocompletion.}\quad For cocomplete $\mathcal{E}$ and $F :
\mathcal{C} \to \mathcal{E}$, $\widetilde F = \operatorname{Lan}_Y F$ is the
essentially unique colimit-preserving extension with $\widetilde F \circ Y
\cong F$. Formula: $\widetilde F(P) = \int^c P(c) \otimes F(c)$.

\medskip
\textbf{Density.}\quad $P \cong \int^c P(c) \otimes \mathcal{C}(-, c)$ for
any $\mathcal{V}$-functor $P : \mathcal{C}^{\mathrm{op}} \to \mathcal{V}$.
$Y$ is dense: $\mathrm{Id}_{[\mathcal{C}^{\mathrm{op}}, \mathcal{V}]} =
\operatorname{Lan}_Y Y$.

\medskip
\textbf{Adjoint as represented hom.}\quad $G : \mathcal{D} \to \mathcal{C}$
has $L \dashv G$ iff $\mathcal{C}(c, G-) \cong \mathcal{D}(L c, -)$ for each $c$.

\medskip
\textbf{Yoneda as unitality in $\mathbf{Prof}$.}\quad $\mathcal{C}(-, -)$ is
the identity 1-cell at $\mathcal{C}$ in $\mathbf{Prof}$: $P \circ \mathcal{C}(-,-)
\cong P$ for all profunctors out of $\mathcal{C}$ (and dually).

\medskip
\noindent\textit{Chapter~29 returns to monoidal categories with the
strictification theorem in full generality (Mac Lane coherence), now in the
2-categorical / bicategorical language of Chapter~25.}
\end{formalrestatement}
